\documentclass[10pt,a4paper]{amsart}
\usepackage[margin=19mm]{geometry}
\usepackage{amsmath,amssymb,mathtools}
\usepackage[scr=rsfs,cal=euler]{mathalfa}
\usepackage{xfrac}
\usepackage[unicode,hidelinks,bookmarks=false]{hyperref}
\newcommand{\ie}{i.e.\ }
\newcommand{\cf}{cf.\ }
\newcommand{\resp}{respectively\ }
\newcommand{\Subsection}{\subsection}
\providecommand{\nobreakdash}{\nobreak\hbox{-}}
\setlength{\emergencystretch}{2em}
\multlinegap0pt
\usepackage{mathtools}
\usepackage{amssymb}
\usepackage[shortlabels]{enumitem}
\usepackage{braket}
\newtheorem{lemma}{Lemma}
\newcommand{\mf}[1]{\mathfrak{#1}}
\newcommand{\rk}{\operatorname{rk}}
\title{Irreducibility of polarized automorphic Galois representations in infinitely many dimensions}
\author{Zachary Feng, Dmitri Whitmore}
\date{Source: arXiv:2507.22631v2 (CC BY 4.0)}
\begin{document}
\maketitle

\noindent\emph{Source and licence.} Irreducibility of polarized automorphic Galois representations in infinitely many dimensions. Version arXiv:2507.22631v2 (CC BY 4.0), distributed by arXiv under the Creative Commons Attribution 4.0 International licence, http://creativecommons.org/licenses/by/4.0/, as recorded in the arXiv licence metadata for this exact version and shown on the abstract page https://arxiv.org/abs/2507.22631v2. Full licence notice: \texttt{licenses/arxiv-2507.22631-CC-BY-4.0.txt}.\par\smallskip

\noindent\textit{Editorial scope.} Counted unit: the article's lemma consequence (source0.tex lines 1020--1023). The article's complete original proof of it is delivered with it (source0.tex lines 1024--1043, one \texttt{proof} environment of 20 lines). No mathematics is written, repaired or re-derived: the statement and the proof are the article's own lines, in the article's own notation, and no retained line is substituted or re-flowed.  No further source-local context is needed: every cross-reference of every retained line resolves inside the two delivered spans.  Nothing was omitted from any retained span, so the package carries no omission record.  No retained line contains a citation, so the wrapper carries no bibliography and no bibliography entry.  No retained line contains a figure environment, an image-inclusion command or drawing code, so no image is delivered and the package declares no asset dependency.  Editorial formatting only, in addition: an AMS \texttt{amsart} preamble that restates the article's own declarations and macros used by the retained lines, namely the environments \texttt{lemma} on separate counters, together with the standard packages those lines need.  The retained environments print in this document as Lemma 1 in order of appearance, whereas the article prints the same items with the numbering of its own full document.  The complete native source of the article as distributed in the arXiv e-print for this exact version is bundled unchanged as \texttt{sources/182431/source0.tex}.
\begin{lemma} \label{lemma: Goursat consequence}
Let $\mf{h} = \prod_{i=1}^k \mf{h}_i$ be a semisimple Lie algebra and let $\mf{f} \subset \mf{h}$ be a semisimple subalgebra such that the induced map $\mf{f} \to \mf{h}_i$ is surjective for every $1 \leq i \leq k$.
If there is an equality of ranks $\rk \mf{f} = \rk \mf{h}$, then there is an equality of Lie algebras $\mf{f} = \mf{h}$.
\end{lemma}

\begin{proof}
We proceed by induction on $k$, first proving the lemma in the case where $k = 2$ (since the case $k=1$ is obvious).
A version of Goursat's lemma for Lie algebras states that there exists ideals $I_i \leq \mf{h}_i$ and an isomorphism $\varphi: \mf{h}_1/I_1 \xrightarrow{\sim} \mf{h}_2/I_2$ such that
\[\mf{f} = \{(x,y) \in \mf{h}_1 \times \mf{h}_2: \varphi(x + I_1) = y + I_2\}.\]
Since each $\mf{h}_i$ is semisimple, each ideal $I_i$ is a semisimple subalgebra and direct product factor of $\mf{h}_i$.
There is a short exact sequence
\[
0 \to I_1 \times I_2 \to \mf{f} \to \mf{h}_1/I_1 \to 0,
\]
from which we see that $ \rk \mf{h}_1 +  \rk I_2 =  \rk \mf{f} = \rk \mf{h}_1 + \rk \mf{h}_2 $.
We must therefore have $\mf{h}_2 = I_2$ and the result in this case now easily follows.

Now suppose that $k>2$.
Let $\mf{h}' = \prod_{i=1}^{k-1} \mf{h}_i$ and let $\pi: \mf{h} \to \mf{h}'$ denote the natural projection.
Let $\mf{f}' = \pi(\mf{f})$.
Then the map $\mf{f}' \to \mf{h}_i$ is surjective for every $1 \leq i \leq k-1$ and we have \[ \rk \mf{h}' = \rk \mf{f} - \rk \mf{h}_k \leq \rk \mf{f'} \leq \rk \mf{h'},\]
so equalities must hold throughout.
By induction, we can therefore apply the lemma to $\mf{f}' \subset \mf{h}'$ to deduce that $\mf{f}' = \mf{h}'$.
We see that $\mf{f}$ surjects onto both $\mf{h}'$ and $\mf{h}_k$, so applying the lemma once again, we see that $\mf{f} = \mf{h}$, which completes the proof.
\end{proof}

\end{document}
